\section{The operator basis}\label{app:D.24}

The reduction of \cref{sec:8.1} takes an interaction written for a dark-matter particle and one free fermion and turns it into operators of the whole target. Four steps take it there: the invariants from which the operators are built, the non-relativistic limits of the Dirac bilinears between plane waves, the one-body operators that reproduce them in a many-body target, and the identities that relate the rows on shell. This section then treats nuclei, the spin traces of the dark-matter-side tensor, and the kinematics of absorption. Natural units are used, and $\vect{q}$ is the momentum the dark matter loses (\cref{sec:1.3}).

\subsection{Galilean invariants}

The dark matter goes from $\vect{p}$ to $\vect{p}' = \vect{p} - \vect{q}$ and the target fermion, of mass $m$, from $\vect{k}$ to $\vect{k}' = \vect{k} + \vect{q}$. A Galilean boost by a velocity $\vect{u}$ shifts $\vect{p}$ and $\vect{p}'$ by $\mchi\vect{u}$, and $\vect{k}$ and $\vect{k}'$ by $m\vect{u}$. The momentum transfer $\vect{q}$ is invariant, and so is the difference of the mean velocities,
\begin{equation}
  \vect{v}^\perp = \frac{\vect{p} + \vect{p}'}{2\mchi} - \frac{\vect{k} + \vect{k}'}{2m} = \bar{\vect{v}}_\chi - \vect{v}_{\mathrm{t}},
  \label{eq:D.24-vperp}
\end{equation}
with $\vect{v}_{\mathrm{t}} \equiv (\vect{k} + \vect{k}')/2m$ (used only in this section). Of the four momenta, conservation leaves three and invariance two, $\vect{q}$ and $\vect{v}^\perp$. Exchanging the incoming and outgoing particles, which is what Hermitian conjugation does to a transition operator, sends $\vect{q} \to -\vect{q}$ and leaves $\vect{v}^\perp$ unchanged; since conjugation also sends $\ii \to -\ii$, the combination $\ii\vect{q}$ is unchanged too. Operators built from $\ii\vect{q}$, $\vect{v}^\perp$ and the spins are therefore Hermitian when their coefficients are real~\cite[Sect.~2.2]{fitzpatrick2013effective}.

For free particles, with the dark matter changing its mass by $\Delta_\chi$ when it is inelastic, energy conservation reads $p^2/2\mchi + k^2/2m = p'^2/2\mchi + k'^2/2m + \Delta_\chi$. Since $p^2 - p'^2 = \vect{q}\cdot(\vect{p} + \vect{p}')$ and $k'^2 - k^2 = \vect{q}\cdot(\vect{k} + \vect{k}')$, this is
\begin{equation}
  \vect{q}\cdot\vect{v}^\perp = \Delta_\chi ,
  \label{eq:D.24-free}
\end{equation}
and $\vect{v}^\perp$ is transverse to $\vect{q}$ for elastic scattering on any free target, whatever its initial momentum. In the rest frame of the target, $\vect{k} = 0$, this is $\vect{q}\cdot\vect{v} = q^2/2\mu_{\chi T}$ with $\mu_{\chi T}$ the reduced mass, as in \cref{sec:1.3}.

\subsection{Bilinears between plane waves}

With the spinors \cref{eq:C.15-spinor}, $u = (E + m)^{1/2}(\xi, \boldsymbol{\sigma}\cdot\vect{k}\,\xi/(E + m))$, and the states normalized to one particle per volume, each bilinear is divided by $(2E\cdot2E')^{1/2}$. Write $\vect{S} = \boldsymbol{\sigma}/2$ and $\vect{K} = \vect{k} + \vect{k}'$. To second order in $k/m$, $[(E + m)/2E]^{1/2} = 1 - k^2/8m^2$ and $1/(E + m) = 1/2m$. With $(\boldsymbol{\sigma}\cdot\vect{k}')(\boldsymbol{\sigma}\cdot\vect{k}) = \vect{k}'\cdot\vect{k} + \ii\boldsymbol{\sigma}\cdot(\vect{k}'\times\vect{k})$ and $\vect{k}'\times\vect{k} = \vect{q}\times\vect{K}/2$, the scalar and the time component of the vector are
\begin{equation}
  \begin{aligned}
  \frac{\bar u'u}{(4EE')^{1/2}} &= 1 - \frac{k^2 + k'^2}{8m^2} - \frac{\vect{k}'\cdot\vect{k}}{4m^2} - \frac{\ii\vect{S}\cdot(\vect{q}\times\vect{K})}{4m^2}
  = 1 - \frac{K^2}{8m^2} - \frac{\ii\vect{S}\cdot(\vect{q}\times\vect{K})}{4m^2} , \\
  \frac{\bar u'\gamma^0u}{(4EE')^{1/2}} &= 1 - \frac{k^2 + k'^2}{8m^2} + \frac{\vect{k}'\cdot\vect{k}}{4m^2} + \frac{\ii\vect{S}\cdot(\vect{q}\times\vect{K})}{4m^2}
  = 1 - \frac{q^2}{8m^2} + \frac{\ii\vect{S}\cdot(\vect{q}\times\vect{K})}{4m^2} .
  \end{aligned}
  \label{eq:D.24-S-V0}
\end{equation}
The second line is the vertex \cref{eq:C.15-vertex} found there by the Foldy--Wouthuysen transformation. The two differ at second order: the scalar loses $K^2/8m^2 = v_{\mathrm{t}}^2/2$, the vector time component only $q^2/8m^2$. The other bilinears are needed to first order. With the blocks of \cref{eq:C.15-gamma} and $\sigma^i\sigma^j = \delta^{ij} + \ii\epsilon^{ijl}\sigma^l$,
\begin{equation}
  \frac{\bar u'\gamma^iu}{2m} = \frac{K^i}{2m} - \frac{\ii(\vect{q}\times\vect{S})^i}{m}, \qquad
  \frac{\bar u'\gamma^0\gamma^5u}{2m} = \frac{2\vect{S}\cdot\vect{K}}{2m}, \qquad
  \frac{\bar u'\gamma^i\gamma^5u}{2m} = 2S^i, \qquad
  \frac{\bar u'\ii\gamma^5u}{2m} = -\frac{\ii\vect{q}\cdot\vect{S}}{m} ,
  \label{eq:D.24-bilinears}
\end{equation}
with the two-component spinors understood on either side. For the spatial vector, $\bar u'\gamma^iu = u'^\dagger\alpha_{\mathrm{D}}^iu = \xi'^\dagger[\sigma^i(\boldsymbol{\sigma}\cdot\vect{k}) + (\boldsymbol{\sigma}\cdot\vect{k}')\sigma^i]\xi$ to this order, and the Pauli algebra gives $K^i - \ii\epsilon^{ijl}q^j\sigma^l$. These agree with Ref.~\cite[Table~I]{trickle2022effective}, whose $\vect{q}$ has the same sign. For the dark matter, whose momentum changes by $-\vect{q}$, replace $\vect{q} \to -\vect{q}$ and $\vect{K}/2m \to \bar{\vect{v}}_\chi$.

A massive vector mediator exchanged between two currents adds the spatial components to the potential of \cref{sec:1.4}. The Proca field obeys $(\partial^2 + m_V^2)A'^\mu = y_\chi J^\mu_\chi$ for a source $-y_\chi A'_\mu J^\mu_\chi$, each component as the time component in \cref{eq:C.9-static}. For $\omega \ll q$ its transform is $y_\chi J^\mu_\chi/(q^2 + m_V^2)$, and a second current with coupling $y$ has the energy $y(A'^0J^0 - \vect{A}'\cdot\vect{J})$. The potential is
\begin{equation}
  \tilde V = \frac{y_\chi y}{q^2 + m_V^2}\,\bigl(J^0_\chi J^0 - \vect{J}_\chi\cdot\vect{J}\bigr) ,
  \label{eq:D.24-vector}
\end{equation}
with the currents the bilinears of \cref{eq:D.24-S-V0,eq:D.24-bilinears}. Its leading term is $\mathcal{O}_1$, and the spatial product is the source of the corrections estimated in \cref{we:8-darkphoton}.

\subsection{From one particle to many}

Each operator of \cref{tab:8-operators} has, between plane waves, a matrix element $f(\vect{k}, \vect{k}')\,\delta_{\vect{k}',\vect{k}+\vect{q}}$ with $f$ a polynomial of degree at most two in $\vect{K}$, times at most one spin. The one-body operator with these matrix elements is unique, and for the three monomials it is
\begin{equation}
  1 \to \ee^{\ii\vect{q}\cdot\vect{r}}, \qquad
  \frac{\vect{K}}{2m} \to \frac12\Bigl\{\frac{\vect{p}}{m}, \ee^{\ii\vect{q}\cdot\vect{r}}\Bigr\}, \qquad
  \frac{K^2}{4m^2} \to \frac{1}{4m^2}\bigl\{\vect{p}, \{\vect{p}, \ee^{\ii\vect{q}\cdot\vect{r}}\}\bigr\} .
  \label{eq:D.24-monomials}
\end{equation}
The proof is one line each: $\langle\vect{k}'|\ee^{\ii\vect{q}\cdot\vect{r}}|\vect{k}\rangle = \delta_{\vect{k}',\vect{k}+\vect{q}}$, and a factor $\vect{p}$ on the left of $\ee^{\ii\vect{q}\cdot\vect{r}}$ contributes $\vect{k}'$ while one on the right contributes $\vect{k}$. The double anticommutator is $p^2X + 2\vect{p}\cdot X\vect{p} + Xp^2$ with $X = \ee^{\ii\vect{q}\cdot\vect{r}}$, with element $k'^2 + 2\vect{k}'\cdot\vect{k} + k^2 = K^2$. Writing $K^2 = q^2 + 4\vect{k}\cdot\vect{k}'$ gives the second form of $\mathcal{T}_s$ in \cref{eq:8.1-rows}. The spin $\vect{S}$ commutes with $\vect{r}$ and $\vect{p}$ and simply multiplies. For a target of many particles of species $s$, the operator is the sum over particles, as for any one-body operator (\cref{eq:2.3-onebody}), and in second quantization it is $\sum_{\vect{k}}f(\vect{k}, \vect{k} + \vect{q})\,c^\dagger_{\vect{k}+\vect{q}}c_{\vect{k}}$ with the spin indices of $\vect{S}$ between the operators. In position space, the density row is the transform $n_e(\vect{q}) = \int\dd^3r\,\ee^{\ii\vect{q}\cdot\vect{r}}\,\op{\nd}_e(\vect{r})$ of the density operator of \cref{sec:4.2}, and the other rows are transforms of the corresponding densities. These are the rows of \cref{eq:8.1-rows}.

\subsection{Reduction of the operators}

With \cref{eq:D.24-vperp}, each $\mathcal{O}_i$ is a polynomial in $\vect{v}_{\mathrm{t}}$ and $\vect{S}_e$ with coefficients built from $\bar{\vect{v}}_\chi$, $\vect{S}_\chi$ and $\vect{q}$. Expanding and replacing $1 \to n$, $S^\beta_e \to s^\beta$, $v^\alpha_{\mathrm{t}} \to j^\alpha$, $v^\alpha_{\mathrm{t}}S^\beta_e \to \mathcal{J}^{\alpha\beta}$ and $v_{\mathrm{t}}^2 \to (2/m)\mathcal{T}$ gives \cref{tab:8-operators}. Three cases show every pattern. For $\mathcal{O}_3$,
\begin{equation*}
  \ii\vect{S}_e\cdot\Bigl(\frac{\vect{q}}{m}\times\vect{v}^\perp\Bigr)
  = \frac{\ii}{m}\Bigl[(\vect{q}\times\bar{\vect{v}}_\chi)\cdot\vect{S}_e - \epsilon_{\alpha\beta\gamma}S^\alpha_eq^\beta v^\gamma_{\mathrm{t}}\Bigr]
  \;\to\; \frac{\ii}{m}\Bigl[(\vect{q}\times\bar{\vect{v}}_\chi)\cdot\vect{s} - \epsilon_{\alpha\beta\gamma}q^\beta\mathcal{J}^{\gamma\alpha}\Bigr] ,
\end{equation*}
where the velocity index of $\mathcal{J}$ comes first. For $\mathcal{O}_{12}$, $(\vect{S}_e\times\vect{v}^\perp)^\alpha = \epsilon_{\alpha\beta\gamma}S^\beta_e(\bar v^\gamma_\chi - v^\gamma_{\mathrm{t}})$, so $\vect{S}_\chi\cdot(\vect{S}_e\times\vect{v}^\perp) \to \vect{S}_\chi\cdot(\vect{s}\times\bar{\vect{v}}_\chi) - \epsilon_{\alpha\beta\gamma}S^\alpha_\chi\mathcal{J}^{\gamma\beta}$. For $\mathcal{O}_{15}$ the same cross product is dotted into $\vect{q}$, $(\vect{S}_e\times\vect{v}^\perp)\cdot\vect{q} \to \vect{q}\cdot(\vect{s}\times\bar{\vect{v}}_\chi) - \epsilon_{\alpha\beta\gamma}q^\alpha\mathcal{J}^{\gamma\beta}$, and multiplied by $-\vect{S}_\chi\cdot\vect{q}/m^2$. The operators without $\vect{S}_e$ need only $n$ and $\vect{j}$, those linear in $\vect{S}_e$ without $\vect{v}^\perp$ only $\vect{s}$, and those with both need $\vect{s}$ and $\mathcal{J}$. Operator $\mathcal{O}_2$ is the one case of degree two, \cref{eq:8.1-O2}.

\subsection{Longitudinal identities}

Let the target Hamiltonian consist of the kinetic energies of its particles and of potentials that depend on their positions alone, as $H_{\mathrm{C}}$ of \cref{eq:2.3-HC} does. Then $[H, n_s(\vect{q})] = \vect{q}\cdot\vect{j}_s(\vect{q})$ for each species, by the commutator of \cref{eq:4.3-continuity} with $m_e \to m_s$, and between eigenstates with $E_f - E_i = \omega$,
\begin{equation}
  \langle f|\vect{q}\cdot\vect{j}_s(\vect{q})|i\rangle = \omega\,\langle f|n_s(\vect{q})|i\rangle .
  \label{eq:D.24-cont}
\end{equation}
If, in addition, $H$ commutes with every spin $\vect{S}_j$ of the species, the same commutator multiplied by $S^\beta_j$ gives $[H, s^\beta_s(\vect{q})] = q^\alpha\mathcal{J}^{\alpha\beta}_s(\vect{q})$, and $\langle f|q^\alpha\mathcal{J}^{\alpha\beta}_s|i\rangle = \omega\langle f|s^\beta_s|i\rangle$. Spin--orbit coupling, the hyperfine interaction and a Zeeman term break this second identity and leave the first intact. On the dark-matter side, \cref{eq:1.3-omegaq} with $\bar{\vect{v}}_\chi = \vect{v} - \vect{q}/2\mchi$ gives $\vect{q}\cdot\bar{\vect{v}}_\chi = \omega + \Delta_\chi$ on shell. Every occurrence of $\vect{q}\cdot\vect{v}^\perp$ in an operator therefore becomes, inside a rate,
\begin{equation}
  \vect{q}\cdot\bar{\vect{v}}_\chi\,n_s - \vect{q}\cdot\vect{j}_s \;\to\; (\omega + \Delta_\chi - \omega)\,n_s = \Delta_\chi\,n_s ,
  \label{eq:D.24-long}
\end{equation}
which is \cref{eq:8.1-transverse}, the many-body form of \cref{eq:D.24-free}. For $\mathcal{O}_2$ the identity applies only to the cross term. The square of the longitudinal electron velocity is not fixed by $\omega$ in a bound target, because the potential energy changes along with the kinetic energy. The three rows of \cref{eq:8.1-O2} are then needed as they stand, and the cancellation of the longitudinal parts used in \cref{we:8-O2} is exact only for free electrons.

\subsection{Nuclei}

Write the nucleons of nucleus $I$ at $\vect{X}_I + \vect{y}_j$, with momenta $(m_N/M_\kappa)\vect{P}_I + \vect{\pi}_j$, where $\vect{P}_I$ is the momentum of the nucleus and $\sum_j\vect{\pi}_j = 0$. A nucleon operator $\sum_j c^{(N_j)}\ee^{\ii\vect{q}\cdot(\vect{X}_I + \vect{y}_j)}$, with $N_j$ the type of nucleon $j$, has, between nuclear ground states, the element $\ee^{\ii\vect{q}\cdot\vect{X}_I}\langle0|\sum_jc^{(N_j)}\ee^{\ii\vect{q}\cdot\vect{y}_j}|0\rangle$. For equal proton and neutron distributions this is $\ee^{\ii\vect{q}\cdot\vect{X}_I}[Z_\kappa c^{(p)} + (A_\kappa - Z_\kappa)c^{(n)}]F_N(q)$ with $F_N$ of \cref{eq:2.2-FN}, which is the first form of \cref{eq:8.1-nuclear}. The nucleon velocity $(\vect{k} + \vect{k}')/2m_N$ splits into $(\vect{P}_I + \vect{P}'_I)/2M_\kappa$ and an internal part built from $\vect{\pi}_j$. The internal part has zero ground-state expectation value, and its contribution at first order in $\vect{q}\cdot\vect{y}_j$ involves the orbital and spin--orbit structure of the nucleus~\cite[Eqs.~(26) and (27)]{fitzpatrick2013effective}; dropping it leaves the rows $\vect{j}_\kappa$ and $\mathcal{T}_\kappa$ with the nuclear mass.

For the spin, $\vect{W} = \sum_jc^{(N_j)}\vect{S}_j$ is a vector operator. Within the multiplet of nuclear spin $J_{\mathrm{nuc}}$ the Wigner--Eckart theorem makes every vector operator proportional to $\vect{J}^{\mathrm{nuc}}$, $\langle M'|\vect{W}|M\rangle = w\,\langle M'|\vect{J}^{\mathrm{nuc}}|M\rangle$. Taking $M' = M = J_{\mathrm{nuc}}$ and the $z$ components, $w\,J_{\mathrm{nuc}} = c^{(p)}\langle S_p\rangle + c^{(n)}\langle S_n\rangle$, where $\langle S_p\rangle = \langle J_{\mathrm{nuc}}, J_{\mathrm{nuc}}|\sum_{j\in p}S^z_j|J_{\mathrm{nuc}}, J_{\mathrm{nuc}}\rangle$ and similarly for neutrons. This is the second form of \cref{eq:8.1-nuclear}. Averaged over $M$ and summed over $M'$ and the components, $\sum_\alpha|\langle M'|J^\alpha_{\mathrm{nuc}}|M\rangle|^2$ gives $\Tr(\vect{J}^2_{\mathrm{nuc}})/(2J_{\mathrm{nuc}} + 1) = J_{\mathrm{nuc}}(J_{\mathrm{nuc}} + 1)$, the factor used for $^{29}$Si in \cref{sec:8.1}. Its ratio to the free-nucleon value $\frac34$ reproduces the spin weight $\tfrac43(J_{\mathrm{nuc}} + 1)J_{\mathrm{nuc}}^{-1}[c^{(p)}\langle S_p\rangle + c^{(n)}\langle S_n\rangle]^2$ of \cref{app:E.27}.

\subsection{Spin traces}

The tensor \cref{eq:8.1-D} is Hermitian, since exchanging $a$ and $b$ conjugates the trace. It is positive semidefinite: for any complex numbers $x_a$, $\sum_{ab}x_a^*\mathfrak{D}_{ab}x_b = (2|\vect{S}_\chi| + 1)^{-1}\Tr_\chi(Y^\dagger Y) \ge 0$ with $Y = \sum_bx_b\tilde V_b$. For a dark-matter spin $S$, $\Tr_\chi\vect{S}_\chi = 0$ and $\Tr_\chi(S^\alpha_\chi S^\beta_\chi) = \delta_{\alpha\beta}(2S + 1)S(S + 1)/3$, so the average of $S^\alpha_\chi S^\beta_\chi$ is $\delta_{\alpha\beta}/4$ for $S = \frac12$. Products of three spin components have a trace proportional to $\epsilon_{\alpha\beta\gamma}$ for spin $\frac12$, and they arise only when a coefficient quadratic in $\vect{S}_\chi$ meets a linear one, which none of the operators of \cref{tab:8-operators} produces.

\subsection{Absorption and the energy deposited}

An absorbed boson enters through $\int\dd^3x\,\phi(x)\,\mathcal{B}(\vect{x})$, with $\mathcal{B}$ a bilinear of the target fields and $\phi \propto \ee^{-\ii m_\phi t + \ii\vect{p}_\phi\cdot\vect{x}}$. The spatial integral is the Fourier transform at $\vect{q} = \vect{p}_\phi$, so the plane-wave rules above apply with $\vect{q} = \vect{p}_\phi$. For the axion-like particle, $\op\psi^\dagger\boldsymbol{\sigma}\cdot\overleftrightarrow{\nabla}\op\psi = \op\psi^\dagger\boldsymbol{\sigma}\cdot\nabla\op\psi - (\nabla\op\psi^\dagger)\cdot\boldsymbol{\sigma}\op\psi$ has the element $\boldsymbol{\sigma}\cdot(\ii\vect{k} + \ii\vect{k}') = 4\ii m_e\,\vect{S}\cdot\vect{K}/2m_e$, since $\nabla$ on the annihilated state gives $\ii\vect{k}$ and on the created state $-\ii\vect{k}'$. The second term of \cref{eq:1.4-alp} is then $(\ii g_{aee}/4m_e^2)(\partial_t\phi)\,4\ii m_e\,\mathcal{J}^{\alpha\alpha}_e = -(g_{aee}/m_e)(\partial_t\phi)\,\mathcal{J}^{\alpha\alpha}_e$, and the first is $-(g_{aee}/m_e)(\nabla\phi)\cdot\vect{s}_e$; the interaction energy is minus the Lagrangian, which is \cref{eq:8.1-alp}. For the scalar, $\op\psi^\dagger\overleftrightarrow{\nabla}^2\op\psi$ has the element $-k^2 - k'^2 - 2\vect{k}'\cdot\vect{k} = -K^2$, and $(1/8m_e^2)(-K^2) = -(1/m_e)(K^2/8m_e)$ is $-1/m_e$ times the element $K^2/8m_e$ of $\mathcal{T}_e$. This is the order-$v^2$ term of the scalar bilinear in \cref{eq:D.24-S-V0}, and with $H = -\mathcal{L}$ it gives \cref{eq:8.1-scalar}.

The energy deposited is \cref{eq:8.1-omega}. For a non-relativistic outgoing particle of mass $m_{\mathrm{out}} = \mchi + \Delta_\chi$,
\begin{equation*}
  \omega = \Bigl(\mchi + \frac{p^2}{2\mchi}\Bigr) - \Bigl(\mchi + \Delta_\chi + \frac{|\vect{p} - \vect{q}|^2}{2m_{\mathrm{out}}}\Bigr)
  = \vect{q}\cdot\vect{v} - \frac{q^2}{2\mchi} - \Delta_\chi + \frac{|\vect{p} - \vect{q}|^2}{2}\Bigl(\frac{1}{\mchi} - \frac{1}{m_{\mathrm{out}}}\Bigr) ,
\end{equation*}
where $p^2 - |\vect{p} - \vect{q}|^2 = 2\vect{q}\cdot\vect{p} - q^2$ was used. The first three terms are $\omega_{\vect{q}}$, and $1/\mchi - 1/m_{\mathrm{out}} = \Delta_\chi/(\mchi m_{\mathrm{out}})$, which is \cref{eq:8.1-omega-inel}. The same formula holds for fermionic absorption into a massive $\tilde\nu$ with $m_{\tilde\nu} = \mchi + \Delta_\chi$. For a massless $\tilde\nu$, $E_{\mathrm{out}} = |\vect{p} - \vect{q}|$, which for $v \to 0$ is $q$, and $\omega = \mchi - q$, \cref{eq:8.1-omega-ferm}. With a free nucleus of mass $M$ as the target, $\omega = q^2/2M$ and $q = \mchi - q^2/2M$, so $q = \mchi(1 - \mchi/2M + \dots)$ and the recoil energy is $\mchi^2/2M$ to first order in $\mchi/M$.
